\documentclass[11pt,a4paper]{article}
\usepackage[margin=27mm]{geometry}
\usepackage{amsmath,amssymb,amsthm,fontspec,unicode-math,hyperref,booktabs,longtable}
\setmainfont{STIXGeneral}
\setmathfont{STIX Two Math}
\setmonofont{Menlo}
\usepackage{newunicodechar}
\newunicodechar{∎}{\ensuremath{\blacksquare}}
\hypersetup{hidelinks}
\setlength{\emergencystretch}{2em}
\providecommand{\tightlist}{\setlength{\itemsep}{0pt}\setlength{\parskip}{0pt}}
\newcommand{\gen}[1]{\langle #1\rangle}
\newcommand{\Sym}{\operatorname{Sym}}
\newcommand{\Alt}{\operatorname{Alt}}
\newcommand{\core}{\operatorname{core}}
\newcommand{\F}{\mathbb F}
\newtheorem{theorem}{Theorem}
\newtheorem{lemma}[theorem]{Lemma}
\newtheorem{proposition}[theorem]{Proposition}
\newtheorem{corollary}[theorem]{Corollary}

\title{Forest coordinates and relative generation}
\author{Human Compression project}
\date{}
\begin{document}
\maketitle
\begin{abstract}
We give two common interfaces for families of flag-transitive geometries.
Positive weights on a tree determine all intersections of the associated
symmetric subgroups, while laminar atoms supply one representative for
every partial flag. Residual maximality, in a second interface, determines
every parabolic from independent choices in the panels. These statements
serve the four applications to inductively minimal geometries, extremal
symmetric-group generators, Frattini quotients and disjoint-subset graphs.
All arguments used by several applications are proved here. The
classification inputs and graph reconstruction arguments that are local
to one application remain in that application.
\end{abstract}

\section{Introduction}
A subgroup intersection calculation and a proof that every flag has a
common representative answer different questions. A forest makes the first
calculation transparent: two vertices remain in one component exactly
when their unique path remains. The same forest, rooted and viewed as a
laminar system of cuts, also answers the second question. Pairwise
incidence determines all nesting and disjointness relations; the atoms
of the resulting partitions can then be matched simultaneously.

Theorem UB.F combines these two facts. Unit weights give the geometries
used by Buekenhout and Cara. Contracting one distinguished edge to a
weight-two vertex supplies the flag calculation for the second family
of Cameron and Cara. Equal leaf weights on a star give Garcia-Colin and
Leemans's disjoint-subset geometry. In each case the application supplies
its actual blocks and cuts, rather than appealing only to the shape of
a diagram.

The other shared argument is an induction down the subgroup intersections.
When at least one next intersection is maximal at every stage, an arbitrary
choice outside each panel stabilizer generates every parabolic. It gives
the relative-generation statements used by both Cameron--Cara and
Archer--Cara--Krempa. The latter paper's Frattini ladder and its finite
$p$-group coordinates are proved in the application, since no other
application in this volume uses those steps.

The unit-weight transposition model, the maximal independent generating
families, and the source geometric results retain their original
attribution. The present organization makes their shared proofs explicit;
it does not claim a new classification of those objects. It also does not
identify a proof of subgroup intersections with a proof of flag completion.

\section{Flags, parabolics and coset coordinates}
All type sets are finite. Groups need not be finite unless this is stated.
A typed incidence system has a reflexive symmetric incidence relation,
with distinct elements of the same type never incident. A flag is a
pairwise incident set. A chamber has every type. A geometry is such a
system in which every flag extends to a chamber. It is firm if every
nonchamber flag has at least two chamber extensions. A residue consists of
the elements incident with a fixed flag and having a type absent from
that flag. Residual connectedness means that every residue of rank at
least two has a connected incidence graph. The specified group action
is part of the data.

For subgroups $H_i\leq G$, indexed by $I$, write
\[
 H_J=\bigcap_{j\in J}H_j,\qquad H_\varnothing=G,\qquad B=H_I.
\]
The elements of type $i$ are the left cosets of $H_i$, and two cosets of
different types are incident if they intersect. The action is left
multiplication. Inversion converts this convention to the right cosets
of the sources. We will use the following three conditions:
\begin{description}
\item[(G1)] The indexed intersections $H_J$ are all distinct.
\item[(G2)] If $|I\setminus J|\geq2$, then
$H_J=\gen{H_{J\cup\{k\}}:k\in I\setminus J}$.
\item[(G3)] Every family of cosets of different types with pairwise
nonempty intersections has a common element.
\end{description}

\noindent\textbf{Proposition UB.Coset (the complete coset interface).}
Conditions (G1)--(G3) hold if and only if the coset system is firm,
residually connected and flag-transitive under its given action.
The residue at the standard flag of type $J$ is the coset geometry of
$H_J$ with subgroups $H_{J\cup\{k\}}$, $k\notin J$.

\begin{proof}
Under (G3), a flag has a representative $g$ common to its cosets and is
therefore contained in the chamber $(gH_i:i\in I)$. Translating $g$ to
the identity shows flag transitivity. Elements of type $k$ incident with
the standard flag of type $J$ can, again by (G3), be represented by an
element of $H_J$. Two such representatives give the same coset precisely
when they differ by $H_J\cap H_k$. This proves the residue description.
Condition (G1) gives $H_{I\setminus\{i\}}>B$, so every panel has at least
two elements. Every shorter flag is contained in a panel; this proves
firmness.

For completeness, a coset incidence graph for subgroups $(L_i:i\in S)$
of a group $L$, with $|S|\geq2$, is connected exactly when the $L_i$
generate $L$. Every edge can be represented by a common group element.
Successive representatives along a path differ by elements of the
subgroup corresponding to their shared vertex. Thus a connected path
from the identity cosets to $gL_i$ places $g$ in $\gen{L_j:j\in S}$.
Conversely, for $a\in L_j$ and $i\ne j$, the vertices
$L_i,L_j=aL_j,aL_i$ give a path. Translating and concatenating these
paths for a word in the $L_j$ connects every coset to the standard ones;
the standard vertices are mutually incident. Applying this criterion
to the described residues proves residual connectedness from (G2).

Conversely, flag transitivity takes any pairwise incident family to the
standard flag of its types, so it has a common representative and (G3)
holds. Firmness and panel transitivity give an element
$s_i\in H_{I\setminus\{i\}}\setminus B$ for every $i$. Its intersection
with $H_i$ is $B$, hence $s_i\notin H_i$. If $J\ne K$, choosing a type
in one set difference shows that $s_i$ belongs to one of $H_J,H_K$ and
not to the other. This proves (G1). Residual connectedness, the residue
description and the graph criterion just proved give (G2).
\end{proof}

An action is primitive when it is transitive and has no nontrivial
invariant partition. A geometry with its group action is PRI if the
action on every type is primitive, and WPRI if the action on at least
one type is primitive. The residual versions, RPRI and RWPRI, apply
to each residue of positive rank, including the whole geometry. In coset
coordinates, RWPRI says
\[
\tag{G4}
J\subsetneq I\quad\Longrightarrow\quad
\text{some }H_{J\cup\{k\}}\text{ is maximal in }H_J,
\quad k\notin J.
\]
This is the elementary equivalence between primitivity on a coset space
and maximality of its stabilizer. It refers to the stabilizer's induced
action; a kernel of that action does not change its invariant partitions.

\noindent\textbf{Lemma UB.Flags (residual weak primitivity gives flags).}
A finite-rank geometry with a residually weakly primitive group action
is flag-transitive.
\begin{proof}
Fix a chamber. In the empty-flag residue choose a type on which the group
is primitive, hence transitive. Match the element of that type in any
second chamber to the fixed one. In the residue of the matched flag,
choose a remaining primitive type and use its flag stabilizer to match
the next element. At most $|I|$ steps match the chambers. Since every
flag extends to a chamber, transitivity follows for flags of every
fixed type. The typewise bijection from cosets of stabilizers to
elements preserves incidence: transitivity on pairs of incident
elements supplies a common representative.
\end{proof}

For $B\leq G$, a finite indexed family $(s_i:i\in I)$ is independent
relative to $B$ if
$s_i\notin\gen{B,s_j:j\ne i}$ for every $i$. If it also generates
$G$ together with $B$, it is a relative generating family. For finite
$G$ let $\mu(G,B)$ and $\mu'(G,B)$ be the largest cardinalities of
relative independent generating families and relative independent
families, respectively. Omit $B$ if $B=1$. Independence for a group
acting on another group means independence under generation by all
images of the other family members; it is denoted $\mu'_G(N)$ in the
normal-quotient application.

\section{Relative generation from residual maximality}
For a finite indexed subgroup family \(H_i\leq G\), write \(H_J=\bigcap_{j\in J}H_j\), with \(H_\varnothing=G\) and \(B=H_I\).

\textbf{Theorem UB.Relative (generation from residual maximality).} Suppose all \(H_J\), \(J\subseteq I\), are distinct, and for every \(J\subsetneq I\) there is \(k\notin J\) such that \(H_{J\cup\{k\}}\) is maximal in \(H_J\). Choose \begin{equation}
s_i\in H_{I\setminus\{i\}}\setminus H_i.
\label{eq:02-relative:6-2}
\end{equation} Then \begin{equation}
H_J=\langle B,s_i:i\notin J\rangle\quad(J\subseteq I).
\label{eq:02-relative:6-3}
\end{equation} In particular these elements generate \(G\) irredundantly relative to \(B\), and their relative parabolics satisfy \begin{equation}
\langle B,s_i:i\in S\rangle\cap\langle B,s_i:i\in T\rangle
=\langle B,s_i:i\in S\cap T\rangle.
\label{eq:02-relative:6-4}
\end{equation} The group need not be finite.

\emph{Proof.} Distinctness of \(H_{I\setminus\{i\}}\) and \(B\) makes \eqref{eq:02-relative:6-2} possible, since its intersection with \(H_i\) is \(B\). Every \(s_j\) with \(j\ne i\), and every element of \(B\), belongs to \(H_i\), while \(s_i\) does not. This proves relative independence.

Prove \eqref{eq:02-relative:6-3} by downward induction from \(J=I\). Choose \(k\) as in the hypothesis. By induction, \(H_{J\cup\{k\}}\) is generated by \(B\) and the \(s_i\) with \(i\notin J\cup\{k\}\). The element \(s_k\) belongs to \(H_J\), because \(J\subseteq I\setminus\{k\}\), and lies outside \(H_{J\cup\{k\}}\), because it lies outside \(H_k\). Maximality therefore says that adjoining \(s_k\) generates \(H_J\). This is \eqref{eq:02-relative:6-3}. Finally substitute \(J=I\setminus S\) and \(I\setminus T\) into the elementary equality \(H_J\cap H_{J'}=H_{J\cup J'}\) to obtain \eqref{eq:02-relative:6-4}. \(\square\)

\noindent\textbf{Theorem UB.Generation (the geometric interface).}
For a finite-rank firm geometry with a residually weakly primitive action,
choose one chamber and its type stabilizers $H_i$. Every choice
\[
 s_i\in H_{I\setminus\{i\}}\setminus H_i
\]
is independent relative to $B=H_I$, generates $G$ over $B$, and satisfies
\[
 H_J=\gen{B,s_i:i\notin J}\qquad(J\subseteq I).
\]
Every such geometry is residually connected. For finite $G$ its rank is
at most $\mu(G,B)$. Without residual weak primitivity, a firm
flag-transitive geometry still has a relatively independent family of
panel choices and, for finite $G$, rank at most $\mu'(G,B)$.
\begin{proof}
Lemma UB.Flags supplies flag transitivity and coset coordinates for the
first assertion. Firmness gives distinct indexed intersections by the
panel argument in Proposition UB.Coset. Residual weak primitivity is
exactly (G4), so Theorem UB.Relative applies to the displayed choices.
Its formula implies (G2): if at least two types remain, each of the
remaining generators belongs to an intersection obtained by fixing
another remaining type. Consequently Proposition UB.Coset gives residual
connectedness. The relative generating family has size equal to the
rank, proving the bound by $\mu(G,B)$.

For the last assertion no maximality is needed. Firmness and flag
transitivity again give panel choices. For a fixed $i$, all other choices
and $B$ lie in $H_i$, whereas $s_i$ does not. They are therefore relatively
independent, which gives the bound by $\mu'(G,B)$.
\end{proof}

\section{Weighted forests and complete flags}
Let \(T\) be a finite tree with vertex set \(V\) and edge set \(E\). Give each vertex a finite nonempty block \(B_v\), with the blocks pairwise disjoint, and put \(\Omega=\bigsqcup_vB_v\). For \(F\subseteq E\), let \begin{equation}
K_F=\prod_C S_{\bigcup_{v\in C}B_v},
\label{eq:03-forest:2-1}
\end{equation} where \(C\) runs over the components of the forest \((V,F)\). Root \(T\), and let \(A_e\subseteq\Omega\) be the union of blocks on the side of \(T-e\) away from the root. In the finite symmetric group \(S_\Omega\), set \(H_e=K_{E\setminus\{e\}}=S_{A_e}\times S_{\Omega\setminus A_e}\).

\textbf{Theorem UB.F (forest intersections and flag completion).} The following hold.

\begin{enumerate}
\def\labelenumi{\arabic{enumi}.}
\tightlist
\item
  \(K_F\cap K_D=K_{F\cap D}\). Moreover \(K_F\) is generated by \(K_\varnothing\) together with the groups \(K_{\{e\}}\), \(e\in F\).
\item
  The coset system of the \(H_e\) is a firm, residually connected, flag-transitive geometry. Its type-\(e\) elements identify with the \(|A_e|\)-subsets of \(\Omega\), and distinct types satisfy \begin{equation}
  X_e\mathrel{*}X_f\quad\Longleftrightarrow\quad
  |X_e\cap X_f|=|A_e\cap A_f|.
  \label{eq:03-forest:2-2}
  \end{equation}
\item
  Its Borel subgroup is \(K_\varnothing=\prod_vS_{B_v}\). The residue whose remaining types are \(F\) is the join of the same constructions on the nontrivial components of \((V,F)\). Its acting group is \(K_F\); components with no edges contribute only a kernel on this residue.
\end{enumerate}

\emph{Proof.} A permutation belongs to both groups in (1) exactly when it preserves every intersection of their component blocks. Two vertices lie in a common block of both forests precisely when their unique tree path belongs to both edge sets. Thus the common refinement is the component partition for \(F\cap D\). Transpositions inside the blocks, together with transpositions joining adjacent blocks along retained edges, connect each component; they generate its full symmetric group. This proves (1).

The coset \(gH_e\) maps to \(gA_e\). Two pairs of subsets are conjugate under \(S_\Omega\) exactly when their four Boolean atoms have the same sizes, giving \eqref{eq:03-forest:2-2}. The sets \(A_e\) are laminar: any two are nested or disjoint. Equation \eqref{eq:03-forest:2-2} forces the identical nesting and disjointness relations on any pairwise incident family \((X_e)\). Add \(\Omega\) as a top container. From each selected container subtract its maximal selected proper subcontainers. These differences partition \(\Omega\), and the corresponding differences for the \(A_e\) and \(X_e\) have equal sizes. Bijections between these atoms combine into a single permutation sending every selected \(A_e\) to \(X_e\). This proves (G3) for every partial flag.

Repeated use of (1) gives \(H_J=K_{E\setminus J}\). Different edge sets give different component partitions and different groups, since every block is nonempty and adding an edge allows a transposition between its two sides. Thus (G1) holds. If at least two edge types remain, the groups obtained by deleting each remaining edge collectively contain \(K_\varnothing\) and each remaining one-edge group, proving (G2). The coset criterion gives the geometry assertions. Its residue group and its subgroups factor over the forest components; all elements from different components are incident, which gives the stated join and the kernel convention. ∎

For unit blocks \(B_v=\{v\}\), write \(\tau_e\) for the endpoint transposition of \(e\). The theorem says \begin{equation}
K_F=\langle\tau_e:e\in F\rangle
 =\prod_C S_C,\qquad K_\varnothing=1.
\label{eq:03-forest:2-3}
\end{equation} This is the unweighted model used for {[}BC98{]}. The weighted formulation also supplies the parity construction after an edge contraction and the disjoint-subset construction after a star specialization.

\section{Overgroups of a two-block stabilizer}
\textbf{Lemma UB.Overgroups (two-block overgroups).} Let \(\Omega=A\sqcup C\) be a finite set, with \(a=|A|\), \(c=|C|\) positive, and put \(H=S_A\times S_C\). If \(a\ne c\), or \(a=c=1\), then \(H\) is maximal in \(S_\Omega\). If \(a=c\geq2\), its only proper strict overgroup is the unordered two-block stabilizer \[
W=(S_A\times S_C)\rtimes C_2,
\] which is maximal and satisfies \([W:H]=2\). In all cases the subgroup interval \([H,S_\Omega]\) is a chain. A positive weighted tree has at most one edge whose two sides have equal total weight.

\emph{Proof.} The case \(a=c=1\) is immediate. The group \(H\) contains the transpositions within each of its two blocks. If an overgroup element does not preserve their unordered partition, its conjugates of these transpositions include one crossing the original blocks. Together with the original within-block transpositions the support graph is connected, so these transpositions generate \(S_\Omega\). This remains valid when one block is a singleton: a permutation moving that singleton produces a cross-block conjugate from the other block. Thus a proper overgroup must preserve the unordered partition. That stabilizer equals \(H\) when the sizes differ and has index two over it when they agree. Its only further overgroup is \(S_\Omega\), by the same argument.

If an edge bisects the total weight, the side of any other edge away from it is a proper positive-weight subset of one of its halves. It has strictly less than half the total weight and cannot be a bisecting side. ∎


\section*{Use and attribution}
The forest theorem is used with unit blocks in \emph{Inductively minimal
geometries and their tree coordinates}, with a contracted distinguished
edge in \emph{Extremal symmetric-group markings and their geometries}, and
with equal leaf blocks in \emph{Disjoint-subset geometries and their full
graph automorphisms}. The relative-generation theorem is used in the
extremal application and in \emph{Frattini quotients and primitive coordinate
geometries}. The two-block overgroup lemma is used in the first two
applications. The full common paper belongs to this topic's proof
dependencies; none of its cost is removed by overlap with another topic.

The coset criterion and relative-generation ideas belong to the established
theory of flag-transitive geometries and to the source works below. We
include elementary proofs in order to expose the exact interfaces. The
weighted-forest atom proof is the organizing argument of the present
volume. Its specializations are not new source results. The substantial
Buekenhout--Cara--Dehon and Whiston classification inputs are stated only
where used and are not concealed inside this common theorem.

\begin{thebibliography}{9}
\bibitem{BC98} F. Buekenhout and P. Cara,
\emph{Some properties of inductively minimal geometries},
Bull. Belg. Math. Soc. Simon Stevin \textbf{5} (1998), 213--219.
\bibitem{CC02} P. J. Cameron and P. Cara,
\emph{Independent generating sets and geometries for symmetric groups},
J. Algebra \textbf{258} (2002), 641--650. The source edition used in this
volume is the twelve-page author manuscript.
\bibitem{ACK05} C. Archer, P. Cara and J. Krempa,
\emph{Using the Frattini Subgroup and Independent Generating Sets to Study
RWPRI Geometries}, Beitr. Algebra Geom. \textbf{46} (2005), 169--177.
\bibitem{GL21} N. Garcia-Colin and D. Leemans,
\emph{An infinite family of locally $X$ graphs based on incidence geometries},
Art Discrete Appl. Math. \textbf{4} (2021), \#P3.01. The exact source
edition used here is the ten-page institutional revised author manuscript,
titled \emph{An infinite family of incidence geometries whose incidence
graphs are locally $X$}.
\end{thebibliography}
\end{document}
